\section{Discussion}\label{sec:discussion}

\paragraph*{Guidance for tool builders}
The results suggest how a slicer for quantum programs should be built. The base should be the wire rule,
in which every gate reads and writes all its qubits. It is sound for every input, misses no dependence,
keeps on average only \numLcExcess{} points more than the dependences, and costs milliseconds. Reading a gate as a classical assignment is the pitfall to avoid, and adding
control dependence does not remove it, since most of the dependences the classical PDG misses are
quantum. Entanglement edges should not be added to make a slice safer. They buy no soundness that the
wire rule lacks and cost precision, because a later operation on an entangled partner cannot change the
criterion's state (Proposition~\ref{prop:ent}). Precision beyond the wire slice can come from
commutation, which stayed correct on every random input but fails when the angles change, so a slicer
for variational programs should recompute it after each parameter update. State-based pruning should be
offered, if at all, as a view of one input, labelled as such, and never as the slice of the program.
Finally, a slicer should let the developer name a set of qubits. Many quantum bugs show only in the
correlation between qubits, as in Section~\ref{sec:example}, and by Proposition~\ref{prop:joint} the
joint wire slice is just the union of the single-qubit slices, so the larger criterion costs nothing
extra.

\paragraph*{Noise and hardware}
A slice is computed statically, so noise cannot change which operations it contains, only whether running
it shows the developer what the full program would. We simulated the pairs of at most ten qubits under
depolarizing noise at three error rates. The noisy state of the criterion under the wire slice differs
from that under the full circuit by at most \numNoiseLcCrossExactMax{} in trace distance where the
simulation is exact, which is rounding error, and by at most \numNoiseLcCrossShotMax{} where mid-circuit
measurements force sampling with \numNoiseShots{} shots. This is what Theorem~\ref{thm:sound} leads one
to expect, since depolarizing noise acts only on the qubits of its gate, so a removed gate and its noise
lie outside the criterion's dependences alike. The state-based slice differs from the full circuit in
\numNoiseQCrossExactNonzero{} of \numNoiseQCrossExactN{} exact runs, by up to \numNoiseQCrossExactMax{}.
It removes noisy gates that the full program executes, and errors accumulate with the number of
gates~\cite{preskill2018nisq}, so it can hide noise the developer is trying to diagnose. On IBM hardware, slices of \numHwCircuits{} circuits from an earlier implementation with
entanglement edges matched the full circuits' measured distributions within a total variation distance of
0.05 for \numHwWithin{}. The other two are QAOA circuits, whose distance is zero in noise-free simulation
and is therefore due to the device, which ran the slice and the full circuit on different physical
qubits. A distance of measured distributions is blind to phase, so these runs show only that slicing
does not break execution on hardware.

\paragraph*{The source level}
A developer reads lines, and a line inside a loop creates one operation per iteration. The line
belongs to the slice when any of them does, so a line-level slice can keep code that no operation-level
slice needs. Fig.~\ref{fig:loops} shows this for the ansatz. In its second layer, the rotations of $q_2$
are not dependences of $q_0$, yet lines 12 and 13 are in the slice. Our tool therefore reports the
iterations of each line. Its static interpretation has limits. A circuit rewritten as a DAG or built
from a QASM string is reported but not traced, and a loop over a range computed at run time is
over-approximated. srcML's markup for OpenQASM~\cite{behler2024supporting, behler2025static} would let
the same back end slice OpenQASM programs.
